% =============================================================
%  Ratings affected by the recovery — short-circuit level, cable
%  stresses, repetitive thermal duty, voltage grading. Principles; the
%  measured cable stresses are in pugalur-thrissur.tex.
% =============================================================
\subsection{Ratings Affected by the Recovery}
\label{sec:dcf-ratings}

\subsubsection{Short-circuit level and power flow}
\label{sec:dcf-ratings-scl}

The short-circuit level of the AC networks acts in opposite directions on
the two fault types~[24]: the pole-to-pole fault current rises with the
short-circuit level, whereas the pole-to-ground overvoltage rises as the
short-circuit level falls, because the load rejection of the faulted link
raises the AC voltage more in a weak network. The pre-fault active and
reactive power have no comparable influence on the fault behaviour beyond
this overvoltage effect. The equipment ratings and the protection settings
are fixed against a defined relation of network overvoltage to
short-circuit level, so that the extreme network conditions are covered
without over-rating~[24].

\subsubsection{Cable stresses}
\label{sec:dcf-ratings-cable}

A hybrid DC circuit raises the fault rate, since the overhead line is
exposed, while it is the cable that has to carry the fault currents and the
overvoltages~[24]. The currents through the cable screens and the grounding
system, the overvoltages across the main insulation and the polarity
reversals of a pole-to-pole fault enter the cable rating. The cable
impedance reduces the $\mathrm{d}u/\mathrm{d}t$ of the initial voltage
collapse as seen at its far termination, and the termination impedance in
the converter station damps the response; steep fronts lie outside the
bandwidth of the installed fault recorders and are covered in the design
studies~[24]. The overvoltage of the healthy pole is not uniform along the
cable: travelling waves and the arresters at the cable ends put the peak
away from the terminals, near the middle of the route in the benchmark of
Section~\ref{sec:dcf-faults-transient}~[26], so the insulation is rated
against the profile along the route. The discharge resistor of Section~\ref{sec:dcf-switchgear}
is part of the same consideration: a resistance well above the surge
impedance keeps the discharge of the healthy pole from launching travelling
waves of its own.

\subsubsection{Repetitive recoveries}
\label{sec:dcf-ratings-thermal}

Each fault and recovery is a thermal event. The devices with long thermal
time constants --- arresters, insertion and discharge resistors, and the
star-point unit --- are rated for the number of recovery sequences the
specification requires within a given period, and their thermal state is
monitored by the protection~[24].

\subsubsection{Voltage grading}
\label{sec:dcf-ratings-grading}

The star-point unit also sets the steady-state grading of the two pole
voltages. The unsymmetrical leakage of a DC overhead line, from insulator
and corona currents, is specified as a maximum grading current unbalance;
for considerably longer overhead sections the supplier states that the
present grading devices may not remain an economic solution, and that
alternative devices or a different main-circuit design would be
needed~[24].

\FloatBarrier
